1: % Reference entries, L (AN-11).  Law 1: os/iword/core/edit.c iw_ed_spacing
2: % ([Ln Spacing:] D0/02 at the cursor, no return before it; 8.8 fixed,
3: % iw_spacing_from_hundredths rounds up), iw_spacing_rows (a line takes its
4: % spacing rounded to the nearest whole screen row), face_a.c M_FMTSP and
5: % follow() (the screen scrolls by rows), the status line's Ln (the spacing
6: % counted).  Shots: ../wordperfect51-decomp mint MINT25, MINT26, MINT30.
7: % GATES: test-pubs-iword test-iword-codes test-iword-wp5-roundtrip test-iword-face test-flair-iword-format
8: \Gates{test-pubs-iword test-iword-codes test-iword-wp5-roundtrip test-iword-face test-flair-iword-format}
9: \begin{Feature}{Line Spacing}
10: Line spacing sets how far apart the lines of text are: single spacing,
11: double spacing for a draft that will be marked up, or any spacing between,
12: such as one and a half.  It applies from the place you set it to the end of
13: the document, or to the next Line Spacing code.
14: 
15: To set the line spacing,
16: \begin{WPSteps}
17: \item Move the cursor to where the new spacing is to begin.
18: \item Press \FK{Format}{Shift-F8}, then select \Ch{Line}{1}.
19: \item Select \Ch{Line Spacing}{6}.
20: \item Type the spacing: \Typed{2} for double spacing, \Typed{1.5} for one and
21: a half, \Typed{1} for single spacing.
22: \item Press \K{Enter}.  The Format: Line screen shows the new spacing.
23: \item Press \FK{Exit}{F7} to return to the document.
24: \end{WPSteps}
25: \Documents{wpmenu}{Shift-F8,1,6=Line Spacing}
26: 
27: \BeginReplay{ref-spacing}
28: For example, press \Press{Format}{Shift-F8}, select \Choose{Line}{1} and
29: \Choose{Line Spacing}{6}, type \TypeName{2}, press \PressKey{Enter} and
30: \Press{Exit}{F7}.  Type \TypeRun{Line one of a double-spaced draft.}, press
31: \PressKey{Enter}, type \TypeRun{Line two.} and press \PressKey{Enter}:
32: \WordScreen[rows=5]{ref-spacing}{Double spacing on the editing screen}
33: \EndReplay
34: 
35: \Notes
36: \Topic{Code}
37: Line Spacing puts a \Code{[Ln Spacing:\textit{n}]} code at the cursor, for
38: example \Code{[Ln Spacing:2]}.  InitechWord asks before deleting one on the
39: editing screen (see \InRef{Codes}).
40: 
41: \Topic{Hundredths}
42: You can type a spacing to two decimal places, for example \Typed{1.33}.
43: InitechWord keeps spacings in steps of 1/256 of a line, and stores a typed
44: spacing as the nearest step at or above it.
45: 
46: \Topic{In the Middle of a Line}
47: Line Spacing goes in exactly at the cursor; it does not end the line there.
48: The line the code is on is spaced by the spacing in effect at its end.
49: 
50: \Topic{On the Screen}
51: On the editing screen each line takes its spacing rounded to the nearest
52: whole line: double spacing leaves one blank row under each line, one and a
53: half spacing also leaves one, and 1.33 leaves none.  The status line's Ln
54: counts the exact spacing, so it shows where the line falls on the page.
55: 
56: \Topic{Typing a Spacing}
57: At the spacing prompt you may type digits and a period.  Backspace erases
58: the last character.  Pressing Enter without typing anything, or after typing
59: something that is not a number of lines up to 8, leaves the spacing as it
60: was; Cancel does the same.
61: 
62: \SeeAlso{Format; Margins, Left and Right; Status Line}
63: \end{Feature}
64: 
65: % List Files (F5): bead initech-qvgd (IWORD P3 remainder, 2026-10-20). Law 1:
66: % WordPerfect 5.1 Reference (initech-manuals-ref/wordperfect51/WPDOS51Reference.pdf)
67: % PDF pp. 364-370 (List Files) and 381-384 (Look); the real-WP mint
68: % ../wordperfect51-decomp mint/work/LISTF/mint-results-005.md (session LF14,
69: % the golden of test-iword-listfiles and test-flair-iword-listfiles);
70: % spec/iword/listfiles.h (the locked screen layout, the 80 -> 76 column
71: % adaptation), os/iword/core/listfiles.c, face_a.c.
72: \begin{Feature}{List Files}
73: List Files shows the documents and other files on the data diskette, with
74: the size and date of each, and lets you retrieve a document or look at one
75: without leaving InitechWord.
76: 
77: To list the files,
78: \begin{WPSteps}
79: \item Press \FK{List}{F5}.  The status line shows \Scr{Dir A:\textbackslash*.*}.
80: \item Press \K{Enter} to list that directory.
81: \item The List Files screen shows the directory, its free space and the
82: number of files on the top line, and the files in two columns.  Use the
83: arrow keys to move the highlight.  \K{Enter} on a directory opens it;
84: \Scr{.. Parent} returns to the one above.
85: \item Select \Ch{Retrieve}{1} to bring the highlighted document onto the
86: editing screen, or \Ch{Look}{6} to read it without retrieving it.  Press
87: \K{Esc} to leave Look, and \FK{Exit}{F7} or \K{Esc} to leave the list.
88: \end{WPSteps}
89: \Documents{wpkey}{F5}
90: If a document is already on the editing screen, Retrieve asks \Scr{Retrieve
91: into current document? No (Yes)}; answer \Typed{Y} to add the file at the
92: cursor, or \Typed{N} to keep the screen as it is.
93: 
94: \Notes
95: \Topic{Entries}
96: Each entry prints the file name, its size in bytes, and the date and time it
97: was last saved, as the operating system's directory listing prints them.  The
98: top line's \Scr{Used} figure is the total of the sizes listed; \Scr{Free} is
99: the space left on the diskette.  Directories print \Scr{<Dir>} in place of
100: a size.
101: 
102: \Topic{Menu items not yet in this version}
103: The List Files menu also names Delete, Move/Rename, Print, Short/Long
104: Display, Other Directory, Copy, Find and Name Search.  This version of
105: InitechWord prints those items in grey and does nothing when they are
106: selected; use the operating system's commands for the same tasks (see the
107: \emph{InitechDOS User's Guide}).
108: 
109: \SeeAlso{Exit; Retrieve; Save}
110: \end{Feature}
111: 
